\documentclass[10pt]{beamer}
\usetheme{metropolis}
\usepackage{graphicx}
\usepackage{listings}
\usepackage{booktabs}
\usepackage{xcolor}
\definecolor{codebg}{HTML}{F5F5F5}
\definecolor{codegreen}{HTML}{2E7D32}
\definecolor{codegray}{HTML}{757575}
\definecolor{codeblue}{HTML}{1565C0}
\lstdefinestyle{rstyle}{language=R,backgroundcolor=\color{codebg},basicstyle=\ttfamily\scriptsize,
  keywordstyle=\color{codeblue}\bfseries,stringstyle=\color{codegreen},
  commentstyle=\color{codegray}\itshape,breaklines=true,frame=single,
  rulecolor=\color{codegray!30},showstringspaces=false,tabsize=2,
  morekeywords={library,read_csv,filter,select,mutate,group_by,summarise,count,
    distinct,drop_na,str_remove_all,as.numeric,fct_reorder,fct_lump_n,fct_infreq,
    ggplot,aes,geom_histogram,geom_boxplot,geom_violin,geom_col,geom_point,geom_smooth,
    geom_errorbar,geom_qq,geom_qq_line,stat_summary,
    coord_flip,scale_x_log10,theme_minimal,labs,patchwork,plot_annotation}}
\lstdefinestyle{pythonstyle}{language=Python,backgroundcolor=\color{codebg},basicstyle=\ttfamily\scriptsize,
  keywordstyle=\color{codeblue}\bfseries,stringstyle=\color{codegreen},
  commentstyle=\color{codegray}\itshape,breaklines=true,frame=single,
  rulecolor=\color{codegray!30},showstringspaces=false,tabsize=4,
  morekeywords={import,from,as,pd,np,plt,sns,query,assign,groupby,agg,
    drop_duplicates,dropna,to_numeric,value_counts,nlargest,sort_values,
    GridSpec,add_subplot,errorbar,hist,scatter,violinplot,barh}}
\title{EDA Case Study: Google Play Store}
\subtitle{Workshop 4 --- Module 7}
\author{Prof.Asoc.\ Endri Raco}
\institute{Data Visualization for Data Scientists \\ UNYT Tirana}
\date{2025/26}
\begin{document}

\begin{frame}
  \titlepage
\end{frame}

\begin{frame}{Module Roadmap}
  \begin{enumerate}
    \item Data cleaning audit trail: 10,841 $\rightarrow$ 9,360 rows
    \item Apply the four-phase EDA workflow to a real dataset
    \item Build a nine-panel EDA dashboard
    \item Extract and communicate four key findings
    \item Outlier deep-dive by category
  \end{enumerate}
  \begin{exampleblock}{Learning Outcome}
    You will execute a complete EDA pipeline on the Google Play Store
    dataset, integrating every technique from W04-M01 through M06
    (cleaning, missingness, outliers, correlation) into a single
    coherent analysis.
  \end{exampleblock}
\end{frame}

\section{Phase 1: Clean \& Document}

\begin{frame}{Data Cleaning Audit Trail}
  \begin{center}
    \includegraphics[width=0.78\textwidth]{figures_w04m07/cleaning_audit}
  \end{center}
  \begin{alertblock}{Pro-Tip}
    Every cleaning step must be \textbf{documented and justified}.
    An audit trail (raw $\rightarrow$ dedup $\rightarrow$ type-fix
    $\rightarrow$ NA drop $\rightarrow$ filter) makes your analysis
    reproducible and defensible. Include row counts at each step.
  \end{alertblock}
\end{frame}

\begin{frame}[fragile]{The Cleaning Pipeline (R)}
\begin{lstlisting}[style=rstyle]
apps_clean <- read_csv("googleplaystore.csv") |>
  # Step 1: Remove 1,181 duplicate app entries
  distinct(App, .keep_all = TRUE) |>
  # Step 2: Fix column types
  mutate(
    Reviews  = as.numeric(Reviews),
    Installs = str_remove_all(Installs, "[+,]") |> as.numeric(),
    Price    = str_remove(Price, "\\$") |> as.numeric(),
    log_reviews = log10(pmax(Reviews, 1)),
    log_installs = log10(pmax(Installs, 1))
  ) |>
  # Step 3: Drop missing Rating (293 rows)
  drop_na(Rating) |>
  # Step 4: Keep valid Type
  filter(Type %in% c("Free", "Paid"), Rating >= 1, Rating <= 5) |>
  # Step 5: Factor ordering for ggplot
  mutate(Category = fct_lump_n(Category, 8) |> fct_infreq())
# Result: 9,360 rows x 15 columns (audit: 10841 -> 9660 -> 9367 -> 9360)
\end{lstlisting}
\end{frame}

\begin{frame}[fragile]{The Cleaning Pipeline (Python equivalent)}
\begin{lstlisting}[style=pythonstyle]
apps_clean = (
    pd.read_csv("googleplaystore.csv")
    .drop_duplicates(subset="App", keep="first")          # 10841 -> 9660
    .assign(
        Reviews=lambda d: pd.to_numeric(d["Reviews"], errors="coerce"),
        Installs=lambda d: d["Installs"].str.replace(r"[+,]", "", regex=True)
            .pipe(pd.to_numeric, errors="coerce"),
        Price=lambda d: d["Price"].str.replace("$", "", regex=False)
            .pipe(pd.to_numeric, errors="coerce"))
    .dropna(subset=["Rating"])                             # -> 9367
    .query("Type in ['Free','Paid'] and Rating>=1 and Rating<=5")  # -> 9360
    .assign(
        log_reviews=lambda d: np.log10(d["Reviews"].clip(lower=1)),
        log_installs=lambda d: np.log10(d["Installs"].clip(lower=1)))
)
\end{lstlisting}
\end{frame}

\section{Phase 2: The Nine-Panel Dashboard}

\begin{frame}{Complete EDA Dashboard}
  \begin{center}
    \includegraphics[width=0.95\textwidth]{figures_w04m07/eda_dashboard}
  \end{center}
\end{frame}

\begin{frame}{Reading the Dashboard: Panel by Panel}
  \centering
  {\small
  \begin{tabular}{rlll}
    \toprule
    \textbf{\#} & \textbf{Panel} & \textbf{EDA Question} & \textbf{Technique (Module)} \\
    \midrule
    (a) & Missing \% bar & How much is missing? & M04 (naniar/missingno) \\
    (b) & Histogram + lines & What shape is Rating? & M01 (Tukey 4-plot) \\
    (c) & Violin by Type & Do Free/Paid differ? & W03-M02 (violin) \\
    (d) & Sorted bar & Which categories dominate? & W03-M03 (bar) \\
    (e) & Scatter + trend & Reviews $\rightarrow$ Rating? & M06 (correlation) \\
    (f) & Corr. heatmap & Which vars correlate? & M06 (heatmap) \\
    (g) & Pointrange & Group means differ? & W03-M07 (CI) \\
    (h) & Waffle & Free/Paid split? & W03-M05 (proportion) \\
    (i) & Q-Q plot & Is Rating normal? & M01 (diagnostics) \\
    \bottomrule
  \end{tabular}}
\end{frame}

\begin{frame}[fragile]{Building the Dashboard (R)}
\begin{lstlisting}[style=rstyle]
library(patchwork)
# (a) Missing audit   (b) Histogram       (c) Violin
# (d) Sorted bar      (e) Scatter          (f) Heatmap
# (g) Pointrange      (h) Proportion       (i) Q-Q

# Compose: 3 x 3 grid
dashboard <- (p_miss | p_hist | p_violin) /
             (p_bar  | p_scatter | p_heat) /
             (p_point | p_waffle | p_qq) +
  plot_annotation(
    title = "Google Play Store: Complete EDA Dashboard",
    subtitle = "9,360 apps | 33 categories | Cleaned & Wrangled",
    caption = "Source: Kaggle | UNYT Data Viz Course",
    tag_levels = "a") &
  theme_minimal(base_size = 8)

ggsave("eda_dashboard.pdf", dashboard, width = 16, height = 11)
\end{lstlisting}
\end{frame}

\section{Phase 3: Extract Findings}

\begin{frame}{Four Key Findings}
  \begin{center}
    \includegraphics[width=0.95\textwidth]{figures_w04m07/findings}
  \end{center}
\end{frame}

\begin{frame}{Finding 1: Left-Skewed Ratings (Ceiling Effect)}
  \begin{itemize}
    \item Mean Rating = 4.17, Median = 4.30 (median $>$ mean
          $\Rightarrow$ left skew)
    \item Most apps cluster at 4.0--4.5; very few below 3.0
    \item Q-Q plot confirms departure from normality at both tails
    \item \textbf{Implication}: use \textbf{Spearman} (not Pearson)
          for correlations; \textbf{median} (not mean) for central
          tendency; consider \textbf{non-parametric} tests
  \end{itemize}
\end{frame}

\begin{frame}{Finding 2: Massive Free/Paid Imbalance (93/7\%)}
  \begin{itemize}
    \item 93\% of apps are Free; only 7\% Paid
    \item Any comparison between Free and Paid must account for this
          imbalance: confidence intervals for Paid are much wider
    \item \textbf{Implication}: stratify by Type before comparing
          categories; use \textbf{bootstrap} or \textbf{weighted}
          methods to avoid sample-size artefacts
  \end{itemize}
\end{frame}

\begin{frame}{Finding 3: Reviews $\leftrightarrow$ Rating Is Weak ($\rho = 0.18$)}
  \begin{itemize}
    \item Spearman $\rho = 0.18$: statistically significant (large $n$)
          but practically negligible
    \item The scatter is a cloud, not a line: high-review apps span the
          full 1--5 rating range
    \item Log transform of Reviews is essential (range: 0 to 78M)
    \item \textbf{Implication}: popularity (reviews/installs) does not
          predict quality (rating). Don't use one as a proxy for the
          other.
  \end{itemize}
\end{frame}

\begin{frame}{Finding 4: Category Differences Are Small}
  \begin{itemize}
    \item SOCIAL (4.30) has the highest mean; MEDICAL (3.95) the lowest
    \item But 95\% CIs overlap for most categories $\Rightarrow$
          differences are likely not significant after multiple testing
    \item \textbf{Implication}: category is a weak predictor of rating;
          within-category variance dominates between-category variance
  \end{itemize}
\end{frame}

\section{Phase 4: Outlier Deep-Dive}

\begin{frame}{Outlier Deep-Dive by Category}
  \begin{center}
    \includegraphics[width=0.82\textwidth]{figures_w04m07/outlier_dive}
  \end{center}
  \begin{exampleblock}{Data Insight}
    All categories show left-tail outliers (low-rated apps). GAME and
    FAMILY have the most outliers by count (largest $n$). The outliers
    are not data errors --- they represent genuinely poorly-rated apps.
    Decision: keep them, as they represent the real distribution.
  \end{exampleblock}
\end{frame}

\section{Synthesis}

\begin{frame}{The EDA Report Structure}
  Every EDA case study follows this structure:
  \centering
  \begin{tabular}{rl}
    \toprule
    \textbf{Section} & \textbf{Content} \\
    \midrule
    1. Data Description & Source, $n$, variables, types \\
    2. Cleaning Audit & Steps, row counts, justifications \\
    3. Missingness Audit & Matrix, \%, mechanism hypothesis \\
    4. Univariate EDA & Distributions, outliers, normality \\
    5. Bivariate EDA & Scatter, correlation, grouped comparisons \\
    6. Key Findings & 3--5 numbered insights with evidence \\
    7. Limitations & What you can't conclude, next steps \\
    \bottomrule
  \end{tabular}
  \vspace{6pt}
  {\small This structure applies to \textbf{every} dataset you will
  encounter in this course (Netflix, e-Car, CRM, Real Estate) and
  in professional practice.}
\end{frame}

\begin{frame}{Key Takeaways}
  \begin{enumerate}
    \item \textbf{Audit trail}: document every cleaning step with row
          counts (10,841 $\rightarrow$ 9,360).
    \item \textbf{Dashboard first}: a 6--9 panel figure gives a
          helicopter view before diving into individual charts.
    \item \textbf{Findings $\neq$ plots}. Extract numbered insights
          with three columns: finding, evidence, implication.
    \item Google Play Store: ratings are ceiling-skewed, Free dominates
          93\%, Reviews $\not\rightarrow$ Rating, category differences
          are small.
    \item The \textbf{EDA report structure} (7 sections) applies to
          every dataset.
  \end{enumerate}
\end{frame}

\begin{frame}{Essential Readings}
  \begin{itemize}
    \item Wickham, H. \& Grolemund, G. (2023). \emph{R for Data Science},
          2nd ed., Chapter 10 (Exploratory Data Analysis).
    \item Peng, R. (2016). \emph{Exploratory Data Analysis with R}.
          Leanpub. (Free online.)
    \item Google Play Store Dataset: Kaggle.
          \texttt{https://www.kaggle.com/lava18/google-play-store-apps}
  \end{itemize}
  \vspace{6pt}
  \textbf{Next Module}: EDA Case Study --- e-Car Loan Pricing (W04-M08)
\end{frame}

\end{document}
